\section{相关工作}
\label{sec:related}

\paragraph{学习式驾驶场景重建。}
前馈模型用从图像观测中学习的推断替代逐场景重复拟合。STORM 预测高斯与场景流以恢复动态户外场景~\cite{yang2025storm}；DGGT 引入免位姿重建、时间可见性和运动引导的插值~\cite{chen2025dggt}。StreetForward、PointForward 和 Instant NuRec 分别探索时序处理、点对齐表示及面向仿真的前馈重建~\cite{yu2026streetforward,chi2026pointforward,huang2026instantnurec}。我们沿用图像恢复模块提供几何和运动证据的作用，并将输出变量组织为背景与规范资产潜码，在新状态渲染之前完成条件补全。

\paragraph{可组合场景与控制。}
Street Gaussians 将跟踪车辆与静态背景分开建模~\cite{yan2024street}，OmniRe 通过适合各类别的运动模型将这一组织方式扩展到多类参与者~\cite{chen2025omnire}。HUGSIM 基于重建内容提供控制驱动的状态更新与交互仿真~\cite{zhou2024hugsim}。这些方法将持久几何、外观与运动分离。我们在生成表示中保留这种分离，使补全表面继承明确的资产身份及其刚体或关节变换。

\paragraph{生成式重建与补全。}
StreetCrafter 根据三维条件生成视图以指导重建，HorizonForge 将可编辑场景内容与视频扩散结合~\cite{streetcrafter,wang2026horizonforge}。WorldSplat 通过潜模型生成像素对齐高斯，并细化其渲染视频~\cite{zhu2025worldsplat}。GenAssets 学习驾驶资产先验，CompleteSplat 根据图像证据补全高斯场景，JRM 则通过潜空间流匹配融合共享对象的未对齐观测~\cite{yang2025genassets,liao2025complete,wu2026jrm}。LatentWorld 围绕场景实体组织语义占据生成~\cite{park2026latentworld}。DriveCompose 以驾驶日志的恢复结果为条件，在统一查询布局上补全可供 RGB 渲染的背景与规范资产。生成作用于未知潜码，观测证据同时约束采样过程和装配后的渲染结果。
